\documentclass[12pt,a4paper]{article}

\usepackage[T1]{fontenc}
\usepackage{lmodern}
\usepackage[utf8]{inputenc}
\usepackage[shorthands=off,brazilian]{babel}
\usepackage[a4paper,margin=2.5cm]{geometry}
\usepackage{booktabs}
\usepackage{array}
\usepackage{longtable}
\usepackage{ragged2e}
\usepackage{csquotes}
\usepackage{enumitem}
\setlist{nosep,leftmargin=*}
\usepackage{titlesec}
\titlespacing*{\subsection}{0pt}{1.6ex plus .4ex}{0.6ex}
\usepackage[hidelinks]{hyperref}
\usepackage[backend=biber,style=abnt,useprefix=true]{biblatex}

\addbibresource{referencias.bib}

\setlength{\parindent}{1.25cm}
\newcolumntype{L}[1]{>{\RaggedRight\arraybackslash}p{#1}}

% Link do projeto no Overleaf (ou Prism). O enunciado exige que ele apareça no PDF.
\newcommand{\linkprojeto}{\url{https://www.overleaf.com/read/rpywrfftkqyv\#70ad2e}}

\begin{document}

\hypersetup{pageanchor=false}
\begin{titlepage}
\centering
{\large Pontifícia Universidade Católica de Minas Gerais\\
Instituto de Ciências Exatas e Informática\\
Ciência de Dados e Inteligência Artificial\par}
\vspace{4cm}
{\LARGE\bfseries Fichamento de artigos científicos\par}
\vspace{1cm}
% Tema do grupo (Ficha, seção 1.2) e recorte usado para escolher e comparar os artigos (etapa 1).
{\large Tema do grupo: desigualdade geográfica no acesso a modelos de inteligência
artificial generativa comercial\par}
\vspace{0.5cm}
{Recorte: como a literatura empírica mede a desigualdade, entre lugares e organizações,
no acesso à inteligência artificial e na adoção de tecnologias de fronteira\par}
\vspace{3cm}
{\large Pedro Dias Soares\\Eduardo Velloso Landeiro\par}
\vfill
{Projeto em Ciência de Dados IV, Prof. Rodrigo Richard Gomes\\
2\textsuperscript{o} semestre de 2026\par}
\vspace{1cm}
{\small Link do projeto LaTeX: \linkprojeto\par}
\end{titlepage}
\hypersetup{pageanchor=true}

\section*{Recorte e seleção dos artigos}

O tema do grupo é a desigualdade geográfica no acesso a modelos de inteligência artificial
generativa comercial. Para comparar artigos, estreitamos o tema para uma pergunta: como a
literatura empírica mede a desigualdade no acesso à IA e na adoção de tecnologias de fronteira,
entre organizações, entre empresas de um país e entre países.

A busca foi feita na OpenAlex, em título e resumo, em 11/09/2026, e está registrada no
Exercício~1. Pedro Dias Soares encontrou McElheran \emph{et al.} (2024), pela expressão \emph{AI
divide}, e Knez (2023), pela combinação \emph{cross-country}, \emph{technology adoption} e
\emph{diffusion}. Eduardo Velloso Landeiro encontrou Ahmed e Wahed (2020), pela expressão
\emph{compute divide}. O grupo é uma dupla, por isso Pedro ficou com dois artigos; cada um fichou o
que encontrou.

O conjunto segue a sugestão do enunciado. Ahmed e Wahed (2020) é o trabalho de referência: é o mais
antigo dos três e fixou o conceito de \emph{compute divide}. McElheran \emph{et al.} (2024) é o
trabalho recente, o primeiro retrato representativo da adoção de IA pelas empresas de um país
inteiro. Knez (2023) é a alternativa, porque explica a desigualdade por preços relativos
(salários), e não por acesso a insumos ou por características das organizações. Os três PDFs
acompanham este documento.

\clearpage
\section{Fichamento 1: Ahmed e Wahed (2020)}

\noindent\textbf{Autoria do fichamento:} Eduardo Velloso Landeiro

\subsection*{1. Referência bibliográfica}
\noindent\fullcite{ahmed2020}

\subsection*{2. Tema do artigo}
A concentração da pesquisa em inteligência artificial (IA) em poucas organizações depois que o
aprendizado profundo ganhou força, e o papel do acesso desigual a poder computacional
(\emph{compute}) nessa concentração. O artigo compara grandes empresas de tecnologia,
universidades de elite e universidades de menor prestígio pela presença nas principais
conferências da área.

\subsection*{3. Problema investigado}
A suspeita de que a IA moderna, por depender de \emph{hardware} caro, virou um campo de quem tem e
quem não tem se apoiava sobretudo em relatos e reportagens (pp.~2--3). Faltava testar de forma
sistemática se a participação na pesquisa de ponta divergiu entre tipos de organização depois de
2012 e se o acesso a computação explica isso. Os autores fazem três perguntas (p.~3): se a pesquisa
ficou mais concentrada, quem são os principais contribuintes e o que aconteceu com quem já atuava
na área.

\subsection*{4. Objetivo do trabalho}
Estimar o efeito causal da ascensão do aprendizado profundo, marcada pelo concurso ImageNet de
2012, sobre a participação de empresas e de universidades de diferentes faixas do \emph{ranking} QS
nas conferências de IA, e mostrar que o acesso a computação é o mecanismo.

\subsection*{5. Hipótese ou pressuposto}
A hipótese é explícita (pp.~4 e 9--10). Partindo de uma função de produção de conhecimento que
depende de conhecimento, habilidades, materiais, equipamento e esforço (eq.~1, p.~9), os autores
preveem que equipamento mais caro reduz a produção de quem não tem acesso a ele: ``we should
observe a divergence between the two groups'' (p.~10), isto é, grandes empresas e universidades de
elite de um lado e as demais universidades do outro. O pressuposto de identificação também é
explícito (pp.~11--12): o sucesso das GPUs no ImageNet 2012 foi inesperado e por isso funciona como
choque exógeno, ligado ao aumento do \emph{compute} disponível e não ao comportamento anterior das
organizações.

\subsection*{6. Base de dados ou corpus}
Artigos de 2000 a 2019 coletados no Scopus, das conferências mais seletivas de computação segundo
o csrankings.org (pp.~12--15). De 63 conferências e 175.491 artigos, ficaram 171.394 artigos de 57
conferências, depois de excluir anos com menos de 25 artigos e a conferência NAACL. As 10
conferências de IA (AAAI, IJCAI, NeurIPS, ICML, KDD, ACL, EMNLP, CVPR, ECCV e ICCV) formam o grupo
tratado e as outras 47, o controle. As afiliações foram classificadas pela Fortune Global 500 de
2018 (46 empresas de tecnologia), pelos \emph{rankings} QS e US News de 2018 e pelo QS de 2011, além
das listas de instituições historicamente negras (HBCU) e hispânicas (HACU). O painel tem 903
observações de conferência por ano (Tabela~1).

\subsection*{7. Metodologia ou abordagem}
O desenho é quase-experimental, com controle sintético generalizado (pp.~10--11). Para cada
conferência de IA, o método usa as conferências de fora da IA para estimar quantos artigos o grupo
teria publicado se 2012 não tivesse acontecido. O efeito (ATT, efeito médio sobre os tratados) é a
diferença média entre o observado e esse contrafactual depois de 2012. O modelo (eq.~2) tem fatores
latentes, efeitos fixos de conferência e de ano e o total de artigos da conferência como controle;
o número de fatores sai de validação cruzada e os erros-padrão, de 2.000 \emph{bootstraps} por
conferência. O desfecho é o número de artigos com ao menos um autor do grupo.

A robustez é testada com placebos de área (Engenharia de Software, Fig.~5) e de data (choque
fictício em 2009, Fig.~6), contagem ponderada pelas afiliações (Tabela~5), \emph{rankings}
alternativos (Tabelas~6 e~7) e um estimador de completamento de matrizes (p.~29). Para o mecanismo,
há uma regressão de efeitos fixos que cruza o choque com uma série anual de \emph{compute} estimada
pela OpenAI (eq.~3, Tabela~8) e uma análise TF-IDF dos resumos, pontuação que mede quanto um termo
é característico dos textos de um grupo (Fig.~8).

\subsection*{8. Principais resultados}
Em relação ao contrafactual, o conjunto das empresas ganhou 46,63 artigos por conferência por ano
(Tabela~2). As 46 grandes empresas de tecnologia ganharam 44,50 (erro-padrão 2,208), contra uma
média de 23 artigos por conferência (Tabela~1). As grandes empresas de fora da tecnologia não
tiveram efeito ($-$0,910), e sem as dez maiores empresas de tecnologia o efeito cai para 19,94. A
Fig.~4 mostra o efeito subindo até cerca de 180 artigos em 2019. Os artigos só de empresas
cresceram cerca de 8 por conferência por ano, e os feitos com universidades, cerca de 42,
concentrados nas universidades do topo do QS (p.~22).

Entre as universidades (Tabela~4), as QS 1--50 ganharam 40,18 (4,49). As faixas 51--100 e
101--200 não tiveram efeito significativo. As 201--300 perderam 14,34 (5,14), e as 301--500
perderam 5,57 (1,54), cerca de 25\% abaixo do contrafactual. As HBCU não tiveram efeito. Os
placebos não mostram efeito (Figs.~5 e~6), e os resultados se mantêm com a contagem ponderada e com
os \emph{rankings} alternativos (Tabelas~5 a~7). Em 20 anos, as HBCU somam 260 artigos e as HACU,
2913; só a Microsoft tem 3302 (pp.~16--17). No mecanismo, a interação entre \emph{compute} e choque é
positiva para empresas e elite e negativa para as universidades QS 101 a 500 (Tabela~8), e as
empresas aparecem mais associadas a termos de aprendizado profundo (Fig.~8).

\subsection*{9. Conclusões dos autores}
A IA está sendo moldada por poucos atores, sobretudo grandes empresas de tecnologia e
universidades de elite, e centenas de universidades de médio e baixo escalão foram deslocadas das
principais conferências (pp.~38--39). Os autores atribuem isso ao acesso desigual a computação, que
chamam de \emph{compute divide}, e dizem ser o primeiro estudo a mostrar que a necessidade de
equipamento especializado pode \emph{desdemocratizar} um campo científico. Defendem intervenção
pública, como uma nuvem nacional de pesquisa e a abertura de dados públicos (p.~38), e afirmam que
o resultado contraria a literatura segundo a qual as empresas vêm recuando da pesquisa pública
(pp.~34--35).

\subsection*{10. Limitações do estudo}
Os autores reconhecem (p.~34) que o estudo cobre só conferências de topo e que o efeito para
empresas nos anos finais pode estar superestimado, se a entrada na IA mudou a participação delas
nas conferências de controle.

Percebi outras. A variável de \emph{compute} da Tabela~8 não é medida por organização: é uma série
anual única, a do maior treinamento conhecido no ano, igual para todos os grupos. A nota~10 (p.~30)
admite que ela ``does not necessarily mean that all organizations need or have access to that
level of computing power'' (p.~30). A escala dessa série também não fecha: o texto diz que ela foi
usada em log (p.~30), mas a Tabela~1 a descreve em escala original, com desvio-padrão de 41.700,14,
e é esse desvio que a p.~31 usa para interpretar os coeficientes.

O controle de total de artigos da conferência (Tabelas~2 e~4) é afetado pelo próprio choque, já
que as conferências de IA cresceram depois de 2012, e isso pode distorcer o efeito. O desfecho é uma
contagem (p.~14), enquanto as Figs.~1 e~2 mostram participações; com conferências de IA maiores,
parte do ganho absoluto das empresas pode ser só crescimento.

Texto e tabelas divergem em vários pontos. A p.~23 fala em 8 artigos a menos para QS 201--300, e a
Tabela~4 mostra $-$14,34. A p.~31 diz que QS 51--100 teve efeito positivo significativo, e a
Tabela~8 mostra $-$0,0003 sem significância. Na mesma p.~31, QS 101--200 aparece primeiro sem
impacto e depois perdendo espaço, e a Tabela~8 pula a coluna~(5). A robustez com completamento de
matrizes fica ``available upon request'' (p.~29), e os outros placebos de área aparecem só na
nota~8 (p.~25). Por fim, o artigo é um \emph{preprint} do arXiv sem revisão por pares na versão de
2020, compara tipos de organização e não países, e as listas HBCU e HACU (p.~16) só valem para os
Estados Unidos.

\subsection*{11. Contribuições}
É a primeira medida sistemática, com desenho contrafactual e placebos, da concentração da pesquisa
em IA depois de 2012. Separa o ganho das grandes empresas de tecnologia do ganho de outras empresas
e das universidades por faixa de \emph{ranking}, e popularizou a expressão \emph{compute divide}.
Para o nosso tema, mostra que o acesso desigual à IA começa antes do produto comercial, já na
produção do conhecimento.

\subsection*{12. Avaliação crítica}
A base é grande e cobre 20 anos (Tabela~1), e o desenho é mais crível que uma comparação simples de
antes e depois: o ajuste antes de 2012 é bom (Fig.~3), há intervalos de confiança (Fig.~4), dois
placebos (Figs.~5 e~6) e três variações de medida e de \emph{ranking} (Tabelas~5 a~7).

\textit{(a) Auditoria da figura principal.} A afirmação central, de que as grandes empresas ganharam
espaço por causa do choque de 2012, está na Fig.~3 (p.~21) e na Tabela~2 (p.~19). Na Fig.~3 o eixo
vertical parte de perto de zero, mas sem o rótulo 0, e não há faixa de incerteza, que só aparece na
Fig.~4. Cada ponto é a média de 10 conferências tratadas, e o ATT médio de 44,50 da Tabela~2 esconde
que quase todo o efeito está nos três últimos anos, quando a série observada vai de cerca de 95
para cerca de 250. O estimador de completamento de matrizes é citado (p.~29), mas não tabulado, e a
participação percentual usada nas Figs.~1 e~2 não aparece na estimação. A Fig.~8 tem outro
problema: não sustenta duas frases da p.~33. No termo \emph{computational cost}, a barra das grandes
empresas é a mais alta e a das QS 301--500, a mais baixa, o contrário do que o texto diz. Em
\emph{support vector machine} e \emph{bayesian}, quem lidera é a elite (QS 1--50), não as
universidades fora dela. Eu esperaria ver o efeito também em participação percentual, com faixa de
confiança, para separar concentração de simples crescimento das conferências de IA.

\textit{(b) Exame de uma afirmação do resumo.} O resumo diz que a divergência entre grandes empresas
e universidades fora da elite ``is driven by access to computing power or compute'' (p.~1). As
únicas peças que tratam disso são a Tabela~8 (p.~30) e a Fig.~8 (p.~34). Na Tabela~8, a série de
\emph{compute} é igual para todos os grupos, então a interação com o choque funciona como uma
tendência pós-2012 das conferências de IA e capta qualquer coisa que tenha crescido no período. A
Fig.~8 mostra sobre o que cada grupo escreve, numa única conferência, e não que \emph{hardware} cada
um tem. Os próprios autores listam três vantagens das empresas que cresceram juntas depois de 2012,
talento, dados proprietários e \emph{compute} (pp.~7--8), e nenhuma análise separa a última das
outras duas. Classificação: ganho de origem não identificada. A divergência está documentada
(Tabelas~2 e~4); a atribuição ao \emph{compute}, e não ao talento ou aos dados, não está.

\clearpage
\section{Fichamento 2: McElheran et al. (2024)}

\noindent\textbf{Autoria do fichamento:} Pedro Dias Soares

\subsection*{1. Referência bibliográfica}
\noindent\fullcite{mcelheran2024}

\subsection*{2. Tema do artigo}
A adoção inicial de inteligência artificial (IA) pelas empresas dos Estados Unidos: quem adotou,
que tecnologias adotou e onde as empresas adotantes ficam, com base no primeiro levantamento
oficial e representativo sobre o assunto.

\subsection*{3. Problema investigado}
Faltava evidência sistemática e representativa sobre o uso efetivo de IA nas empresas (Introdução,
p.~375). As medidas disponíveis estimavam o \emph{potencial} da IA de forma indireta, por descrições
de tarefas, anúncios de vagas e patentes (p.~377), e as pesquisas privadas que mediam uso tinham
amostras enviesadas para empresas grandes e de capital aberto (p.~382). Sem o \emph{denominador}
certo, que inclui as muitas empresas pequenas e jovens, não dá para saber quão difundida a IA está
nem onde seus efeitos devem aparecer primeiro.

\subsection*{4. Objetivo do trabalho}
Descrever a difusão inicial de cinco tecnologias associadas à IA na economia americana e
identificar as características de empresas, donos, estratégias e locais que se associam ao uso.
O objetivo é descritivo, nas palavras dos autores: ``Our main objective is to document'' (p.~377)
fatos estilizados, ``rather than explain'' (p.~377).

\subsection*{5. Hipótese ou pressuposto}
O artigo não testa uma hipótese única, mas declara expectativas explícitas na Seção~2
(pp.~378--380): donos com mais escolaridade e experiência tenderiam a liderar empresas mais
intensivas em tecnologia; donos com motivação de \emph{estilo de vida} adotariam menos; e capital
de risco e capital inicial alto sinalizariam potencial de crescimento ligado à IA (todas na
p.~379).

Há também um pressuposto que reconstruí: quem adota cedo antecipa quem vai adotar depois, de modo
que um retrato de 2017 indica a trajetória futura da IA. Ele aparece quando os autores falam dos
adotantes iniciais ``and, by extension, the most likely later'' (p.~377) e quando dizem que o
caminho da IA dependerá da sua dispersão geográfica inicial (p.~380). Não é testado, porque os
dados são de um único ano.

\subsection*{6. Base de dados ou corpus}
O levantamento é o \emph{Annual Business Survey} (ABS) de 2018, do Census Bureau com a NSF, com
referência ao ano de 2017. Foi enviado a 850.000 empresas, teve taxa de resposta de 69\% e cerca de
583.000 respondentes (pp.~376 e~380). As respostas foram ligadas à \emph{Longitudinal Business
Database} (LBD), base administrativa com emprego e receita ao longo da vida das empresas. A
amostra-base tem 447.000 empresas, que representam cerca de 4.050.000 com os pesos; a amostra com
dono e receita tem 209.000; e a de \emph{startups} de 0 a 5 anos, 75.000 (Tabela~1, p.~381). O mapa
usa cerca de 30.000 \emph{startups} das áreas metropolitanas com mais de um milhão de habitantes
(Fig.~5, p.~390).

Conta como uso de IA a resposta positiva para ao menos uma de cinco tecnologias: veículos guiados
automaticamente, aprendizado de máquina, visão computacional, processamento de linguagem natural e
reconhecimento de voz (Tabela~2, p.~382). Qualquer intensidade vale, inclusive \emph{menos de 5\% da
produção}.

\subsection*{7. Metodologia ou abordagem}
O estudo é observacional, descritivo e de corte transversal. A primeira parte (Seção~3,
pp.~380--390) traz estatísticas descritivas ponderadas: taxas de uso por tecnologia, porte, setor,
idade, características do dono e área metropolitana, com pesos de empresa, para representar a
população, e pesos de emprego, para medir a exposição dos trabalhadores (Figs.~1--5 e Tabelas~3
e~5).

A segunda parte (Seção~4, pp.~390--396) usa modelos de probabilidade linear por mínimos quadrados.
O uso de IA, uma variável 0/1, é regredido sobre porte, idade, digitalização e nuvem (Tabela~4) e,
para as \emph{startups}, sobre crescimento inicial da receita, características do dono,
financiamento, inovação e porte da área metropolitana (Tabela~7). Todos os modelos têm controles de
estado por indústria de seis dígitos, e cada coeficiente se lê como variação, em pontos percentuais,
da probabilidade de usar IA. Antes disso, os autores validam os marcadores de crescimento
regredindo o crescimento médio da receita nos três primeiros anos sobre eles (Tabelas~6a e~6b,
eq.~1).

\subsection*{8. Principais resultados}
Só 5,8\% das empresas usavam alguma das cinco tecnologias, e 1,1\% estavam testando; ponderado por
emprego, o uso sobe para 18,2\% (Fig.~1, p.~383). O aprendizado de máquina lidera, com 2,9\%, e os
veículos guiados ficam em 0,8\%. A maioria das empresas com mais de 5.000 empregados usa IA
(Fig.~2a). Indústria e informação ficam em torno de 12\%, e construção e varejo, em torno de 4\%
(Fig.~2b, p.~384). Usar nuvem se associa a 5 pontos percentuais a mais na probabilidade de usar IA
(0,050), e ter informação digitalizada, a 1,9 ponto (0,019) (Tabela~4, p.~386). Entre as empresas
que usam robôs, 57,5\% também usam IA (p.~385).

Entre as \emph{startups} (Tabela~7, coluna~5, pp.~395--396, $R^2 = 0{,}230$), os maiores
coeficientes são capital de risco (+8,6 pontos; 0,0862), proteção de propriedade intelectual
considerada \emph{muito importante} (+6,1), inovação de processo (+5,4) e patentes (+3,5). Inovação
de processo e propriedade intelectual, juntas, somam 11,5 pontos, quase o dobro da taxa média de 6\%
(p.~393). Um aumento de 10 pontos no crescimento inicial da receita eleva a probabilidade em 0,13
ponto na coluna~1 e em 0,06 na coluna~5.

Na geografia, os polos ficam acima da média no Vale do Silício e em metrópoles do Sul e do Oeste,
como Nashville, San Antonio, Las Vegas e Tampa (Fig.~5, pp.~389--390); com peso de emprego, entram
Louisville, Columbus, Austin e Atlanta (Fig.~A2). Estar numa área metropolitana grande eleva o uso
em 0,74 ponto (coluna~4, 0,0074), mas o efeito deixa de ser significativo quando entram as
características da empresa (coluna~5, 0,0042).

\subsection*{9. Conclusões dos autores}
O uso inicial de IA era baixo na média, mas concentrado em empresas grandes, inovadoras e com
marcadores de alto crescimento, e em alguns lugares (Seção~5, pp.~396--399). Dados e nuvem são
tecnologias habilitadoras centrais, e por isso os autores esperam um impacto econômico maior do que
a taxa média sugere. A conclusão diz que um possível \emph{AI divide} entre tipos de empresa e entre
lugares ``is visible'' (p.~399) e precisa ser avaliado, e reconhece que estabelecer causalidade
está ``beyond the reach of this first'' (p.~399) corte transversal de dados.

\subsection*{10. Limitações do estudo}
Os autores reconhecem (pp.~391 e~399) que a análise detalhada só vale para empresas jovens com dono
identificado, o que exclui as de capital aberto; que há viés de sobrevivência, porque empresas que
adotaram IA e fecharam antes de 2018 não aparecem; que não há causalidade; e que cerca de 80\% da
variação no uso de IA fica sem explicação (p.~399).

Percebi outras. A medida é frouxa e autodeclarada: conta como uso qualquer resposta positiva,
inclusive em \emph{menos de 5\%} da produção (Tabela~2 e nota da Tabela~7). A nota~2 (p.~400)
admite que os resultados refletem sobretudo \emph{machine learning}, e a definição desse termo na
Tabela~2 é ampla o bastante para abranger usos triviais de modelos preditivos. Os dados também são
antigos para o nosso tema: a referência é 2017, antes da IA generativa comercial, e a taxa de 5,8\%
(Fig.~1) não diz nada sobre o acesso a modelos como os de hoje.

O ganho de ajuste das características organizacionais é pequeno. A frase de que elas
``contributes significantly to the explanatory power'' (p.~399) se apoia num $R^2$ que passa de
0,209 para 0,230 (Tabela~7), ou 2,1 pontos absolutos, apresentados como aumento de ``about 10\%''
(p.~393). A maior parte do $R^2$ vem dos controles de estado por indústria, e o artigo não mostra um
modelo só com esses controles. As descritivas são ponderadas (p.~381), mas as notas das Tabelas~4
e~7 não dizem se as regressões usam os pesos; marco como \emph{[não informado]}. Por fim, a
geografia é só uma comparação entre metrópoles de um país: a Fig.~5 cobre áreas com mais de um
milhão de habitantes (nota~21, p.~401), sem intervalo de confiança por área. Para a pergunta do
nosso grupo, que compara países, o artigo serve de método, não de resultado.

\subsection*{11. Contribuições}
A contribuição principal é de medição. É a primeira estimativa representativa do uso de IA pelas
empresas de um país inteiro, com o denominador certo e ligada a dados administrativos de emprego e
receita (pp.~377 e~380--382). Mostra também que a taxa por empresa e a taxa por emprego contam
histórias diferentes (5,8\% contra 18,2\%, Fig.~1). Para o nosso projeto, ficam duas lições: a
escolha da unidade e do peso muda o tamanho da desigualdade medida, e a desigualdade geográfica
bruta pode ser composicional, já que o efeito do porte da área metropolitana some quando entram as
características das empresas (Tabela~7, colunas~4 e~5).

\subsection*{12. Avaliação crítica}
A amostra é muito grande e representativa, com alta taxa de resposta (69\%, p.~376), e vem de uma
fonte oficial ligada a registros administrativos, o que evita o viés das pesquisas privadas
(p.~382). Os autores são cuidadosos com causalidade e repetem que os achados são correlações
(pp.~391, 396 e~399 e nota~28). Também separam o peso por empresa do peso por emprego, o que é raro.

\textit{(a) Auditoria da tabela principal.} A afirmação central sobre quem adota se apoia na
Tabela~7 (pp.~395--396). A tabela traz erros-padrão, mas a soma de coeficientes que o texto destaca
(11,5 pontos para inovação de processo mais propriedade intelectual, p.~393) aparece sem erro-padrão
próprio. Por ser tabela, a questão do eixo vertical não se aplica. As colunas~1--5 têm 75.000
observações, número arredondado por sigilo, mas os coeficientes de capital de risco dependem de uma
fração pequena delas, porque só 0,9\% das \emph{startups} têm esse financiamento (Tabela~5, p.~388).
Falta um modelo só com os controles de estado por indústria, que mostraria quanto do $R^2$ de 0,209
a 0,230 vem das variáveis de interesse. Também falta a taxa de uso prevista para cada perfil de
empresa, que o texto só apresenta em termos relativos, como ``a 17\% increase at the mean''
(p.~393). Eu esperaria ver a coluna de base só com controles e as somas de coeficientes destacadas
no texto com seus intervalos de confiança.

\textit{(b) Exame de uma afirmação do resumo.} O resumo afirma que a adoção ``was far from evenly
distributed'' (p.~375) e que um punhado de cidades e polos emergentes liderou a adoção pelas
\emph{startups}. A peça principal é a Fig.~5 (p.~390), e ela mostra posições, não taxas: a cor de
cada área é o percentil da sua taxa de uso entre as grandes metrópoles. Uma escala de percentis vai
sempre de 0 a 100, mesmo que as taxas sejam quase iguais, e a figura não mostra a taxa em si nem a
incerteza de cada área. As outras peças apontam diferenças pequenas. O próprio texto diz que as
áreas menores ficam ``not too far below'' (p.~390) as maiores; na Tabela~5, 55,1\% das
\emph{startups} que usam IA estão em grandes metrópoles, contra 49,9\% das que não usam; e na
coluna~5 da Tabela~7 o efeito do porte da área deixa de ser significativo. Classificação: conclusão
mais ampla que o experimento permite. O artigo sustenta uma ordenação das metrópoles, não que a
desigualdade entre elas seja grande nem que seja distinguível de ruído.

\clearpage
\section{Fichamento 3: Knez (2023)}

\noindent\textbf{Autoria do fichamento:} Pedro Dias Soares

\subsection*{1. Referência bibliográfica}
\noindent\fullcite{knez2023}

\subsection*{2. Tema do artigo}
Como tecnologias se espalham entre países com níveis de desenvolvimento muito diferentes. O artigo
quer saber por que a mesma tecnologia chega cedo a alguns países e tarde, ou quase nunca, a outros,
e liga isso ao custo relativo de implantá-la em cada lugar, sobretudo ao salário.

\subsection*{3. Problema investigado}
Há duas literaturas que não se conversam (Seção~2, pp.~1173--1176). Os estudos de curva de adoção,
de Griliches (1957) a Comin e Hobijn (2010), reproduzem a forma logística da adoção de cada
tecnologia, mas tratam a economia como fechada e homogênea; explicam o atraso dentro de um país, não
entre países. As teorias macroeconômicas do desenvolvimento desigual (hiato tecnológico, restrição
de balanço de pagamentos, neoschumpeterianas) explicam divergência e armadilhas de baixo
crescimento, mas não geram curvas por tecnologia e tratam a capacidade de absorção como exógena,
dada por instituições e escolaridade. Falta um modelo em que as curvas de adoção de cada país saiam
das próprias condições econômicas desiguais.

\subsection*{4. Objetivo do trabalho}
Propor um modelo macroeconômico dinâmico de difusão em que a distância entre países não é
geográfica, e sim o custo relativo de implantar a tecnologia, e calibrá-lo com dezenas de
tecnologias e mais de uma centena de países para testar se o salário relativo aproxima essa
distância.

\subsection*{5. Hipótese ou pressuposto}
A hipótese é explícita (Introdução, p.~1172, retomada na p.~1180): ``relative wage level is one of
the most important endogenous socioeconomic determinants of the relative cost of technology
implementation'' (p.~1172). Ou seja, quanto menor o salário de um país em relação à média mundial,
mais cara fica, em termos relativos, a adoção, e mais tarde a tecnologia chega.

Três pressupostos, que reconstruí a partir do texto, sustentam a hipótese sem serem testados: o
capital tem preço único mundial e o salário é nacional (pp.~1180--1181, eq.~3.9); as tecnologias
novas poupam sobretudo trabalho, o que torna decisiva a razão entre salário e preço do capital
(p.~1180); e toda tecnologia começa nos países de menor custo relativo, com a condição inicial
concentrada em $x=0$ (p.~1179, eqs.~3.5--3.7).

\subsection*{6. Base de dados ou corpus}
Dois bancos públicos (Seção~4, pp.~1181--1182). O CHAT (\emph{Cross-country Historical Adoption of
Technology}) é um painel desbalanceado com o uso de tecnologias em mais de 161 países ao longo de
cerca de 200 anos, com medidas extensivas (parcela da população com acesso) ou intensivas
(quantidade em uso). A Penn World Table fornece o salário médio de cada país. O autor selecionou 36
tecnologias de agricultura, infraestrutura, siderurgia, informação, saúde, automóveis e turismo, e
excluiu algumas, como aço e irrigação, por refletirem especialização ou geografia mais do que
difusão (nota~4, p.~1181). As estimações usam 32.552 observações (Tabelas~1 e~2). O ano de invenção
de cada tecnologia é a primeira observação não ausente no banco (p.~1182).

\subsection*{7. Metodologia ou abordagem}
O artigo é teórico, com calibração empírica. A variável é a densidade tecnológica relativa $\Phi$:
o uso per capita de uma tecnologia num país dividido pela média mundial no mesmo ano, o que permite
comparar tecnologias medidas em unidades diferentes (p.~1177). Por analogia com a lei de Fick, o
fluxo de tecnologia vai das regiões mais densas às menos densas, na proporção do gradiente ao longo
do eixo $x$ de custo relativo de implementação (eq.~3.1), o que leva à equação clássica de difusão
(eq.~3.4). Com toda a tecnologia concentrada em $x=0$ na invenção, a solução é uma gaussiana que se
alarga com o tempo (eq.~3.8, derivada no Apêndice~A).

O custo relativo é aproximado por $x = \alpha \cdot (\textrm{salário médio mundial}/\textrm{salário
do país})$ (eq.~3.10), e restam dois parâmetros: a constante de difusão $D$ e a escala $\alpha$. A
estimação é feita por mínimos quadrados não lineares (Tabela~1) e por um modelo não linear de
efeitos mistos (eq.~4.2, Tabela~2), em que $D$ e $\alpha$ ganham um desvio aleatório por tecnologia;
na prática, cada tecnologia tem a própria velocidade em torno de uma média comum. A robustez vem de
2000 reamostragens de 10 tecnologias (Apêndice~B, Figs.~3 e~4, Tabela~3). As Figs.~1 e~2 são
simulações do modelo com os parâmetros estimados.

\subsection*{8. Principais resultados}
Por mínimos quadrados não lineares, $D = 0{,}00202$ e $\alpha = 0{,}329$, ambos com $p < 0{,}001$,
e $R^2 = 0{,}474$ (Tabela~1): dois parâmetros explicam cerca de 50\% da variação da densidade
relativa. Com efeitos mistos, $D = 0{,}00318$ e $\alpha = 0{,}477$, e as variâncias dos efeitos
aleatórios são 0,105 para $\alpha$ e 1,26e-05 para $D$ (Tabela~2). O autor lê essas variâncias como
sinal de que a difusão é ``relatively universal'' (p.~1183) entre as tecnologias. Na reamostragem,
a estimativa de $D$ por mínimos quadrados é bimodal, o que indica diferença entre tecnologias
antigas e modernas, e fica mais estável com efeitos mistos (Figs.~3 e~4, Tabela~3).

Na simulação (Fig.~2, p.~1184), a densidade relativa de um país com o dobro do salário médio sobe
quase logo após a invenção. Com 50\% do salário médio, ela fica perto de zero por mais de 10 anos;
com 25\%, por muitas décadas. Países perto da subsistência ficam ``practically technologically
blocked'' (p.~1185). A adoção de cada país continua logística em termos absolutos, mas o
deslocamento e a inclinação da curva dependem do salário relativo (p.~1185).

\subsection*{9. Conclusões dos autores}
A difusão de tecnologia é um mecanismo endógeno que perpetua o desenvolvimento desigual: o salário
relativo determina a adoção, e a adoção desigual reproduz a diferença de salários, num processo
``dynamically acting in both directions'' (p.~1189). O autor lista quatro contribuições
(pp.~1188--1189): tratar a difusão por uma equação de difusão de fato, trocar o espaço geográfico
pelo de custo relativo, derivar curvas de adoção relativas por país e compatibilizar o resultado
com as teorias evolucionárias e de causação cumulativa. Na discussão, afirma que o modelo gera de
forma endógena a armadilha da renda média (p.~1187) e que fechar o hiato exige política industrial
ativa de transferência de tecnologia, comparando a China com os novos membros da União Europeia no
Leste Europeu (p.~1188).

\subsection*{10. Limitações do estudo}
O autor reconhece (p.~1189) que o modelo abstrai as particularidades de cada tecnologia e de cada
país e não explica a curva de uma tecnologia ou de um país específico.

Percebi outras. O canal salarial não fica isolado: na Seção~4.1 (p.~1186) o autor admite que o
salário também sinaliza capital humano, instituições e desenvolvimento geral, mas nenhuma variável
concorrente (PIB per capita, escolaridade, índice institucional) entra na eq.~4.1. Por isso o ajuste
da Tabela~1 não separa custo de implantação de desenvolvimento em sentido amplo. Também falta um
modelo de comparação: o $R^2$ de 0,474 (Tabela~1) não é confrontado com nenhuma alternativa, como
uma logística por país ou o modelo de Comin e Hobijn citado na p.~1174, e o modelo misto (Tabela~2)
não traz medida de ajuste.

A inferência parece otimista. As 32.552 observações formam um painel de país, tecnologia e ano
tratado como independente, e estatísticas $t$ de 71 e 93 (Tabela~1) sugerem erros-padrão
subestimados, porque não há agrupamento por país. A data de invenção é a primeira observação
disponível no CHAT (p.~1182), o que mede o início da cobertura do banco e não a invenção, e desloca
o eixo do tempo de tecnologias antigas, como ferrovia e telégrafo. A amostra está mal descrita: a
lista da p.~1182 soma 37 tecnologias (4+7+2+8+12+2+2), não 36, e o grupo de sete de infraestrutura
nomeia só três. Os cálculos só estão disponíveis a pedido (Declarações, p.~1193).

Para o nosso tema há ainda um limite de alcance. O CHAT termina antes da IA generativa, e todas as
tecnologias da p.~1182 têm instalação física local. Num modelo de IA oferecido pela nuvem, o custo
de chegar a um país depende da decisão do provedor de ativar a região, não da troca entre salário e
máquina pressuposta na eq.~3.9.

\subsection*{11. Contribuições}
A contribuição é conceitual. O artigo transforma a metáfora da difusão tecnológica numa equação de
difusão de verdade e troca o espaço geográfico por um espaço de custo relativo. Com dois
parâmetros, reproduz o padrão de adoção tardia e lenta nos países de baixa renda (Fig.~2). Para o
nosso projeto, oferece um vocabulário, o atraso de chegada como função de uma distância de custo
entre países, e uma predição testável: se o mecanismo valer para IA comercial, países de renda
relativa menor deveriam receber novos modelos mais tarde, com atraso que cresce de forma não linear
com a distância de renda.

\subsection*{12. Avaliação crítica}
O artigo saiu numa revista de economia evolucionária com revisão por pares, em acesso aberto, usa
dados públicos, mostra a derivação completa (Apêndice~A) e responde à heterogeneidade entre
tecnologias com efeitos mistos (Seção~4, Fig.~3). É o único dos três que trata da dimensão entre
países, que é a do nosso tema.

\textit{(a) Auditoria da figura principal.} A figura que sustenta a tese central, adoção tardia e
mais lenta nos países de salário baixo, é a Fig.~2 (p.~1185). O eixo vertical começa em zero. Não
há dispersão: as curvas usam as estimativas pontuais $D = 0{,}00318$ e $\alpha = 0{,}477$, sem a
faixa dos intervalos de 95\% que a própria Tabela~2 fornece. Nenhuma observação aparece, porque as
curvas são simulação e as 32.552 observações das Tabelas~1 e~2 não são sobrepostas a elas. O
painel~(b) depende ainda de uma relação linear entre salário e renda nacional bruta per capita de
2021 que é suposta, não estimada (legenda da Fig.~2). As curvas logísticas da literatura citada nas
pp.~1174 e~1183 não aparecem, e falta uma medida de ajuste do modelo misto (Tabela~2).

As tabelas trazem mais dois problemas. A variância de 1,26e-05 do efeito aleatório de $D$
(Tabela~2) equivale a um desvio-padrão de cerca de 0,0035, maior que o próprio $D = 0{,}00318$, o
que contradiz a frase ``the differences in diffusion constants are relatively small'' (p.~1183). E
na Tabela~3 os rótulos das colunas parecem trocados: 0,0018 aparece sob $\alpha$ e 0,377 sob $D$, o
inverso da ordem de grandeza da Tabela~1. Eu esperaria ver, para ao menos uma tecnologia, as
densidades observadas por faixa de salário sobrepostas às curvas da Fig.~2, com a faixa de
incerteza dos parâmetros.

\textit{(b) Exame de uma afirmação da conclusão.} A conclusão (p.~1189) afirma que a difusão
perpetua o desenvolvimento desigual nos dois sentidos, do salário para a adoção e da adoção para o
salário, e que isso explica por que países de renda média continuam de renda média. O sentido
inverso não está no modelo: na eq.~4.1 o salário entra como dado observado e exógeno, e nenhuma
equação ou estimativa faz a adoção alterar o salário. A persistência também não decorre da equação.
Pela eq.~3.8, a razão entre as densidades de dois países é $\exp[-(x_1^2-x_2^2)/(4Dt)]$, que tende a 1
quando $t$ cresce; o modelo prevê convergência lenta, não estados estacionários múltiplos. Com os
parâmetros da Tabela~2, um país com 25\% do salário médio tem, cem anos após a invenção, cerca de
6\% da densidade de um país com o dobro do salário médio, e em mil anos, cerca de 75\%. O bloqueio
por décadas está sustentado; a perpetuação e a armadilha da renda média (p.~1187) não estão.
Classificação: conclusão mais ampla que o experimento permite. A parte sobre o efeito da adoção nos
salários é especulação apresentada como explicação.

\clearpage
\section{Quadro comparativo}

{\footnotesize
\begin{longtable}{@{}L{2.1cm}L{4.2cm}L{4.2cm}L{4.2cm}@{}}
\toprule
& \textbf{Ahmed e Wahed (2020)} & \textbf{McElheran et al. (2024)} & \textbf{Knez (2023)} \\
\midrule
\endhead
\bottomrule
\endlastfoot
\textbf{Problema} &
A pesquisa de ponta em IA se concentrou em poucas organizações depois de 2012? Por falta de
acesso a computação? &
Quem usa IA nas empresas dos EUA, que tecnologias e onde? Faltava uma medida representativa. &
Por que tecnologias chegam tarde aos países pobres? Os modelos de adoção ignoram o
desenvolvimento desigual. \\ \addlinespace
\textbf{Dados} &
171.394 artigos de 57 conferências de computação, 2000--2019 (Scopus). Afiliações classificadas
por Fortune 500 e QS. &
ABS 2018 (ano-base 2017) ligado à LBD: 447 mil empresas (4 milhões ponderadas); 75 mil
\emph{startups}. &
CHAT (36 tecnologias, mais de 161 países, cerca de 200 anos) e Penn World Table (salários);
32.552 observações. \\ \addlinespace
\textbf{Método} &
Controle sintético generalizado, com o ImageNet 2012 como choque. Regressão de efeitos fixos com
série de \emph{compute}. TF-IDF dos resumos. &
Descritivas ponderadas e modelos de probabilidade linear (MQO) com controles de estado por
indústria. Corte transversal. &
Modelo de difusão (equação de Fick) no espaço de custo relativo. Calibração por MQ não lineares e
efeitos mistos. \\ \addlinespace
\textbf{Baseline} &
Contrafactual sintético a partir de 47 conferências de fora da IA. Placebos de área e de data. &
Empresas que não usam IA. Não há modelo só com controles. &
Nenhum modelo concorrente. A logística é citada, mas não estimada. \\ \addlinespace
\textbf{Métrica} &
ATT em artigos por conferência por ano, com IC por \emph{bootstrap}. &
Probabilidade de uso (pontos percentuais), taxas ponderadas por empresa e por emprego, $R^2$. &
Parâmetros $D$ e $\alpha$, $R^2$ (só no MQ não linear), curvas simuladas. \\ \addlinespace
\textbf{Resultado} &
Grandes empresas de tecnologia +44,5 e QS 1--50 +40,2; QS 201--300 $-$14,3 e QS 301--500 $-$5,6
(Tabelas~2 e~4). &
5,8\% das empresas usam IA (18,2\% por emprego). Associação com capital de risco, inovação de
processo e propriedade intelectual. Efeito do local some com controles. &
$R^2 = 0{,}474$. Países com 25\% do salário médio ficam décadas perto de zero na densidade
relativa. \\ \addlinespace
\textbf{Limitação} &
\emph{Compute} não é medido por organização. Talento, dados e \emph{compute} não são separados.
\emph{Preprint}. &
Uma única fotografia de 2017, antes da IA generativa. Uso autodeclarado com limiar baixo. Só os
EUA. &
Salário confunde custo com desenvolvimento. Figuras simuladas, sem dados. Sem código. \\
\end{longtable}
}

\section{Síntese comparativa}

\subsection*{Convergência}
Os três trabalhos concordam que a adoção de tecnologia de fronteira começa concentrada em quem já
tem vantagem e tende a se reforçar. Em \textcite{ahmed2020}, grandes empresas de tecnologia e
universidades de elite ganham espaço nas conferências de IA, enquanto as universidades de médio e
baixo escalão perdem. Em \textcite{mcelheran2024}, a IA aparece sobretudo em empresas grandes,
financiadas por capital de risco e inovadoras, e em alguns polos urbanos. Em \textcite{knez2023},
países de salário alto adotam logo após a invenção e os de salário baixo ficam décadas para trás.
Os três também explicam a concentração por complementos, e não pela tecnologia em si:
\emph{compute}, dados e talento em Ahmed e Wahed; dados digitalizados e nuvem em McElheran
\emph{et al.}; a relação entre salário e preço do capital em Knez. E chegam à mesma preocupação: a
desigualdade inicial pode virar desigualdade persistente.

\subsection*{Divergência}
Os três divergem sobre o mecanismo, e a diferença vem sobretudo de método e de unidade de análise.
Para Ahmed e Wahed, pesa a falta de acesso a um insumo, a computação. Para McElheran \emph{et al.},
pesam características da empresa e do dono, como escolaridade, financiamento e estratégia de
inovação. Para Knez, pesam preços relativos definidos pelo mercado, em especial o salário.

A discordância mais clara é sobre o papel do lugar. Em Knez, a posição do país, dada pelo salário
relativo, é o determinante principal da chegada da tecnologia. Em McElheran \emph{et al.}, a
vantagem das grandes áreas metropolitanas some quando se controlam as características das
empresas (Tabela~7, colunas~4 e~5), o que sugere uma desigualdade geográfica composicional:
empresas inovadoras se juntam nos mesmos lugares. Isso não é necessariamente um desacordo de fato.
Knez usa um único regressor e não testa variáveis concorrentes, McElheran \emph{et al.} controlam
dezenas de características, e as escalas são muito diferentes, metrópoles de um mesmo país, com a
mesma regulação e a mesma moeda, contra países de renda muito distinta. As recomendações também
divergem: Ahmed e Wahed pedem oferta pública de computação, Knez defende política industrial de
transferência de tecnologia e McElheran \emph{et al.} pedem dados melhores.

\subsection*{Comparabilidade}
Os resultados não são comparáveis em magnitude. Os desfechos diferem: artigos por conferência por
ano em relação a um contrafactual, em Ahmed e Wahed; probabilidade de uma empresa declarar uso, em
McElheran \emph{et al.}; densidade de uso relativa à média mundial, em Knez. As unidades e os
períodos não se sobrepõem: conferências entre 2000 e 2019, empresas americanas em 2017 e países ao
longo de dois séculos, com tecnologias como o trator e a internet. As métricas também têm naturezas
distintas, porque o ATT de Ahmed e Wahed é uma diferença de contagens, os coeficientes de McElheran
\emph{et al.} são pontos percentuais e o $R^2$ de Knez mede o ajuste de um modelo. Um ganho de 44
artigos e um aumento de 8,6 pontos percentuais não estão na mesma escala.

O que dá para comparar é o sinal e a forma: nos três casos, quem tem mais recurso chega antes ou
chega mais. Uma comparação quantitativa exigiria um desfecho comum, como o tempo entre o lançamento
de uma tecnologia e seu uso ou disponibilidade em cada unidade, e nenhum dos três mede isso
diretamente. Knez chega mais perto, mas com curvas simuladas e sem dados sobrepostos (Fig.~2).

\subsection*{Lacuna}
Somados, os três deixam sem resposta uma pergunta do lado da oferta: quanto tempo um modelo
comercial de IA generativa leva para ficar disponível em cada país ou região, e do que depende
esse atraso. Ahmed e Wahed medem quem produz pesquisa, não quem consegue usar o produto. McElheran
\emph{et al.} medem a demanda das empresas de um único país, antes da IA generativa. Knez trata de
tecnologias com instalação física local, em que o custo de implantar depende do salário; num
modelo entregue pela nuvem, a chegada a um país depende antes da decisão do provedor de ativar o
serviço na região e da regulação local, e o mecanismo salarial pode não valer. Medir esse atraso
por país e por modelo e compará-lo com renda, infraestrutura de nuvem e regulação é o que o nosso
projeto pretende fazer. A mesma medida serve de teste da predição de Knez e da hipótese
composicional de McElheran \emph{et al.}

\clearpage
\printbibliography[title={Referências}]

\end{document}
